\section{Summary and Looking Ahead}
\label{sec:differential-calculus-summary}

Functions describe dependence between quantities, while limits make controlled approach precise. A derivative is the best local linear approximation to a function. In one dimension it is represented by a slope; in several dimensions it is represented by a Jacobian matrix. Partial derivatives provide coordinate rates, but differentiability is the stronger statement that one linear map controls change in every direction.

For scalar fields, the total differential can be written as $df=\nabla f\cdot d\mathbf r$. The gradient points in the direction of steepest ascent, is normal to regular level sets, and determines every directional derivative. Along a path $\mathbf r(t)$, the chain rule gives $df/dt=\nabla f\cdot\dot{\mathbf r}$.

\begin{booksummary}[Chapter 2 in One Line]
Differential calculus replaces a nonlinear function, near one point, by its unique best first-order linear model.
\end{booksummary}

\begin{roadmapbox}[Bridge to Chapter~\ref{ch:vector-calculus}]
Chapter~\ref{ch:vector-calculus} applies the same local-derivative viewpoint to vector fields. Divergence measures local expansion, curl measures local rotation, and the Laplacian combines second derivatives. The integral theorems of Green, Gauss, and Stokes then connect these local quantities to global flux and circulation.
\end{roadmapbox}
